% introduction.tex
%
% Flow:
%   1. Hook: LLMs as optimization oracles — the opportunity
%   2. Problem: current approaches lack rigor
%   3. Gap analysis: what is missing
%   4. Contributions: anneal framework
%   5. Results preview
%   6. Paper structure

\section{Introduction}
\label{sec:introduction}

% --- 1. Hook ---
% Describe the general opportunity: LLMs can propose improvements to software
% artifacts (prompts, code, configs). Give 2-3 concrete motivating examples.
\TODO{Write the hook paragraph: LLMs as optimization oracles. Cite 2-3 recent
examples of LLM-guided optimization (AlphaEvolve, OPRO, DSPy).}

% --- 2. Problem Statement ---
% Describe how existing approaches treat optimization as ad hoc trial-and-error:
%   - No statistical significance testing between iterations
%   - Greedy hillclimbing without diversity maintenance
%   - No cross-run knowledge accumulation
%   - Expensive stochastic evaluation without adaptive stopping
\TODO{Write the problem paragraph. Key point: existing systems produce results
that cannot be reliably reproduced or compared across runs.}

% --- 3. Gap Analysis ---
% Enumerate the specific gaps:
%   G1: Statistical rigor — no hypothesis testing over optimization runs
%   G2: Structured search — no strategy selection or diversity mechanisms
%   G3: Knowledge accumulation — no memory across optimization sessions
%   G4: Evaluation efficiency — stochastic eval cost not controlled
\TODO{Write the gap analysis paragraph.}

% --- 4. Contributions ---
% We present anneal, making the following contributions:
\TODO{Introduce anneal. Frame as a framework, not a tool.}

The contributions of this paper are:
\begin{itemize}
  \item \textbf{Framework}: \textsc{anneal}, an open-source framework for
    artifact-agnostic LLM-guided optimization with a clean
    artifact--eval--agent abstraction (\cref{sec:system-design}).
  \item \textbf{Enhancements}: \NUM{16} enhancements across four tracks
    (Representation, Search, Knowledge, Evaluation) with design rationale
    (\cref{sec:enhancements}).
  \item \textbf{Evaluation}: A benchmark suite of 5 targets across 3 domains,
    with 10 independent runs per configuration and Wilcoxon signed-rank
    testing with Holm-Bonferroni correction (\cref{sec:evaluation}).
  \item \textbf{Empirical Results}: \TODO{Summarize top-line result here after
    benchmarks are complete.} (\cref{sec:results}).
\end{itemize}

% --- 5. Results Preview ---
% One paragraph summarizing the key experimental finding.
% Fill in after Task 5.4.
\TODO{Write results preview paragraph. Reference specific targets where
treatment outperforms control.}

% --- 6. Paper Structure ---
The remainder of this paper is organized as follows.
\Cref{sec:related-work} surveys related work.
\Cref{sec:system-design} describes the \textsc{anneal} architecture.
\Cref{sec:enhancements} details the four enhancement tracks.
\Cref{sec:evaluation} describes the experimental setup.
\Cref{sec:results} presents empirical results.
\Cref{sec:discussion} discusses limitations and threats to validity.
\Cref{sec:conclusion} concludes.
